\chapter{Systemarchitektur}
\label{ch:architektur}

\section{Technologie-Stack}

Die Umsetzung von WikiPod basiert auf mehreren Komponenten, die jeweils eine klar
abgegrenzte Aufgabe innerhalb der Gesamtarchitektur übernehmen. Für die lokale
Bereitstellung der Wikipedia-Inhalte werden ZIM-Dateien im KIWIX-Format verwendet.
Der Zugriff auf diese Dateien erfolgt in Python über \texttt{libzim}. Die in den
Artikeln enthaltenen HTML-Dokumente werden mit \texttt{BeautifulSoup} verarbeitet,
um Textinhalte und Metadaten zu extrahieren.

Für die semantische Repräsentation der Textabschnitte kommt die Bibliothek
\texttt{sentence-transformers} zum Einsatz. Die erzeugten dichten Vektoren werden
gemeinsam mit den zugehörigen Metadaten in OpenSearch gespeichert und dort für
das k-NN-basierte Retrieval verwendet. Die Auswahl von KIWIX beziehungsweise ZIM
als Datenformat und OpenSearch als Retrieval-System ergibt sich unmittelbar aus
der Aufgabenstellung des Projekts.

Die lokale Antwortgenerierung wird über ein kleines Sprachmodell realisiert.
Dafür unterstützt die Implementierung ein direkt eingebundenes
\texttt{llama.cpp}-Backend; für die lokale Entwicklung kann alternativ Ollama
verwendet werden. Die Kommandozeilenschnittstelle des Projekts ist mit
\texttt{Click} umgesetzt und nutzt \texttt{Rich} für strukturierte
Terminalausgaben. Konfigurationen werden mit \texttt{PyYAML} eingelesen und über
\texttt{Pydantic}-Modelle validiert.

\section{Pipeline-Ueberblick}

Die Systemarchitektur von WikiPod lässt sich in zwei voneinander getrennte
Phasen unterteilen: die Indexierungsphase und die Anfragephase.
Abbildung~\ref{fig:pipeline} zeigt den vollständigen Datenfluss von der
Wikipedia-ZIM-Datei bis zur generierten Antwort. Die Trennung ist insbesondere
für den Einsatz auf ressourcenbeschränkter Hardware relevant, da die
rechenintensive Aufbereitung des Datenbestands unabhängig von der späteren
Beantwortung einzelner Nutzeranfragen durchgeführt werden kann.

In der Indexierungsphase bildet eine lokale Wikipedia-ZIM-Datei die
Datengrundlage. Zunächst werden die für die Dokumentenauswahl benötigten
Metadaten extrahiert und in einem JSONL-Cache zwischengespeichert. Auf dieser
Grundlage werden die Link-Frequenzen bestimmt, die Artikel bewertet und unter
Berücksichtigung des verfügbaren Speicherbudgets ausgewählt. Erst für die
ausgewählten Artikel wird anschließend der vollständige Text geladen. Dieser
wird in überlappende Textabschnitte zerlegt und durch ein Embedding-Modell in
dichte Vektorrepräsentationen überführt. Die resultierenden Chunks werden
zusammen mit ihren Embeddings und Metadaten im OpenSearch-Index gespeichert.

Die Anfragephase greift auf diesen vorbereiteten Index zurück. Eine
Nutzeranfrage wird zunächst mit demselben Embedding-Modell wie die
Dokument-Chunks in einen Vektor überführt. Über eine Ähnlichkeitssuche in
OpenSearch werden anschließend die zur Anfrage passendsten Textabschnitte
ermittelt. Diese bilden gemeinsam mit der ursprünglichen Anfrage und den
Anweisungen an das Sprachmodell den Prompt. Das lokal ausgeführte Sprachmodell
erzeugt daraus schließlich die Antwort. Dadurch bleiben sowohl der
Wikipedia-Datenbestand als auch Retrieval und Antwortgenerierung innerhalb der
lokalen Systemumgebung.

\begin{figure}[H]
    \centering
    \resizebox{\textwidth}{!}{%
    \begin{tikzpicture}[
        node distance=7mm,
        phase/.style={
            draw,
            rounded corners=3mm,
            inner sep=8mm
        },
        box/.style={
            draw,
            rounded corners=2mm,
            align=center,
            text width=5.2cm,
            minimum height=10mm
        },
        store/.style={
            draw,
            cylinder,
            shape border rotate=90,
            aspect=0.28,
            align=center,
            text width=2.6cm,
            minimum height=2.8cm
        },
        arrow/.style={
            -{Latex[length=2.2mm]},
            thick
        }
    ]

    % Indexierungsphase
    \node[box] (zim)
        {Wikipedia-ZIM\\{\small Dump der Artikeldaten}};

    \node[box, below=of zim] (metadata)
        {Metadatenextraktion\\und JSONL-Cache};

    \node[box, below=of metadata] (selection)
        {Link-Frequenz-Map,\\Scoring und Auswahl};

    \node[box, below=of selection] (fulltext)
        {Volltext der ausgewählten\\Artikel laden};

    \node[box, below=of fulltext] (chunking)
        {Chunking\\{\small Text in überlappende Abschnitte}};

    \node[box, below=of chunking] (embedding)
        {Embedding der Chunks\\{\small Vektorrepräsentationen}};

    % Anfragephase
    \node[box, right=8.2cm of zim] (query)
        {Nutzeranfrage\\{\small natürliche Sprache}};

    \node[box, below=of query] (queryembed)
        {Query-Embedding\\{\small Vektorrepräsentation}};

    \node[box, below=of queryembed] (retrieval)
        {Dense Retrieval in OpenSearch\\
        {\small ähnliche Chunks suchen}};

    \node[box, below=of retrieval] (prompt)
        {Prompt-Konstruktion\\
        {\small relevante Chunks + Anweisung}};

    \node[box, below=of prompt] (llm)
        {Lokales Sprachmodell};

    \node[box, below=of llm] (answer)
        {Antwort\\{\small natürliche Sprache}};

    % Zentraler OpenSearch-Vektorindex
    \node[
        store,
        right=10mm of embedding,
        yshift=9mm
    ] (opensearch)
        {OpenSearch-\\
        Vektorindex\\
        {\small Chunk-Embeddings\\und Metadaten}};

    % Phasenrahmen
    \begin{scope}[on background layer]

        \node[
            phase,
            fit=(zim)
                (metadata)
                (selection)
                (fulltext)
                (chunking)
                (embedding),
            label={
                [font=\bfseries\Large, anchor=north]
                north:Indexierungsphase
            }
        ] (indexphase) {};

        \node[
            phase,
            fit=(query)
                (queryembed)
                (retrieval)
                (prompt)
                (llm)
                (answer),
            label={
                [font=\bfseries\Large, anchor=north]
                north:Anfragephase
            }
        ] (queryphase) {};

    \end{scope}

    % Indexierungsfluss
    \draw[arrow] (zim) -- (metadata);
    \draw[arrow] (metadata) -- (selection);
    \draw[arrow] (selection) -- (fulltext);
    \draw[arrow] (fulltext) -- (chunking);
    \draw[arrow] (chunking) -- (embedding);

    \draw[arrow]
        (embedding.east)
        --
        (opensearch.west);

    % Anfragefluss
    \draw[arrow] (query) -- (queryembed);
    \draw[arrow] (queryembed) -- (retrieval);
    \draw[arrow] (retrieval) -- (prompt);
    \draw[arrow] (prompt) -- (llm);
    \draw[arrow] (llm) -- (answer);

    % OpenSearch -> Dense Retrieval
    \draw[arrow]
        (opensearch.north)
        -- ++(0,12mm)
        -| (retrieval.west);

    \end{tikzpicture}%
    }

    \caption{Überblick über die WikiPod-Pipeline mit Indexierungs- und Anfragephase}
    \label{fig:pipeline}
\end{figure}

\section{Projektstruktur und Konfiguration}

Die Implementierung von WikiPod ist als Python-Paket aufgebaut und gliedert die
einzelnen Verarbeitungsschritte der Pipeline in eigenständige Module. Dadurch
bleiben Dokumentenauswahl, Indexierung, Retrieval und Antwortgenerierung
voneinander getrennt, können jedoch über die Kommandozeilenschnittstelle zu
einem durchgängigen Ablauf verbunden werden. Ergänzend enthält das Projekt
separate Verzeichnisse für Konfigurationsdateien, Tests und den vorliegenden
Bericht.

Ein zentraler Bestandteil der Projektstruktur ist die konfigurationsbasierte
Steuerung des Systems. Allgemeine Standardwerte werden in
\texttt{config/default.yaml} definiert und können durch umgebungsspezifische
Konfigurationsdateien überschrieben werden. Dadurch lassen sich beispielsweise
Entwicklungs-, Test- und produktive Indexierungsläufe mit unterschiedlichen
Ressourcenbudgets und Parametern durchführen, ohne dafür Änderungen am
Programmcode vornehmen zu müssen.

Über die Konfiguration werden unter anderem Parameter der Dokumentenauswahl,
des Chunkings, des Embedding-Modells, der OpenSearch-Anbindung, des Retrievals
und der lokalen Antwortgenerierung festgelegt. Die YAML-Konfiguration wird beim
Programmstart eingelesen und anschließend durch Pydantic-Modelle validiert.
Fehlende oder ungültige Einstellungen können dadurch frühzeitig erkannt
werden.

Dieser Ansatz vermeidet zugleich, dass experimentell bestimmte Parameter als
fest codierte Werte über verschiedene Module verteilt werden. Dies ist für
WikiPod insbesondere deshalb relevant, weil sich geeignete Einstellungen
zwischen Entwicklungsumgebung und Zielhardware deutlich unterscheiden können.
Die für einen konkreten Lauf verwendeten Parameter bleiben auf diese Weise
nachvollziehbar und können für Evaluationen reproduziert und dokumentiert
werden.